\section{Adaptación de dominio y post training}
\label{sec:marco-post-training}
Durante el preentrenamiento, un modelo de lenguaje desarrolla capacidades generales a partir de grandes corpus, mientras que la adaptación de dominio busca orientar esas capacidades hacia el vocabulario, los conocimientos y las tareas de un ámbito específico. En este contexto, el proceso de post training puede combinar demostraciones supervisadas con la evaluación de respuestas generadas por el propio modelo, de modo que cada etapa aporta una señal de aprendizaje diferente~\cite{gururangan2020dontstop,ouyang2022instructgpt}.

Dentro de las estrategias de adaptación, el preentrenamiento adaptativo al dominio (DAPT) continúa el aprendizaje sobre textos especializados, mientras que la adaptación al corpus de una tarea (TAPT) utiliza textos asociados a esa tarea~\cite{gururangan2020dontstop}. Por su parte, el SFT utiliza demostraciones explícitas de cómo responder, de manera que incorporar documentos del área y entrenar con pares de consulta y respuesta no constituye el mismo procedimiento. En consecuencia, su alcance depende de la calidad y la cobertura de las referencias, cuyos errores u omisiones también pueden aprenderse.

La elección de la señal depende también de la tarea. En problemas con resultados verificables puede utilizarse una regla que compruebe la solución, mientras que en dominios abiertos algunos aspectos admiten comprobaciones automáticas y otros requieren interpretar el contenido y el contexto~\cite{deepseek2025r1,arora2025healthbench}. Esta diferencia motiva estudiar cómo construir recompensas a partir de criterios explícitos, después de establecer los fundamentos del ajuste supervisado y del aprendizaje por refuerzo.
